\documentclass[12pt]{article}
\usepackage[T1]{fontenc}
\usepackage[utf8]{inputenc}
\usepackage{lmodern}
\usepackage{textcomp}
\usepackage{enumitem}
\usepackage{geometry}
\geometry{margin=1in}
\begin{document}
\begin{center}
Answer Key - Sociology - Graduate Level Assessment 5 \
\small (Answers are ordered exactly as in the question paper.)
\end{center}

\section*{Multiple Choice}
\begin{enumerate}[label=\arabic*.]
\item \textbf{Answer:} D
\\ \textit{Options:} A) grants every member equal status; B) does not have any official secrets in its government; C) has permissive attitudes towards sexual behaviour; D) allows people to move between levels of the hierarchy
\\ \textit{Note:} An 'open' society is characterized by social mobility, which enables individuals to change their socio-economic status and move between different levels of the social hierarchy, reflecting a more egalitarian structure.
\item \textbf{Answer:} C
\\ \textit{Options:} A) the findings are amenable to statistical analysis; B) it is conducted over a period of several years; C) it uncovers rich, detailed accounts from an insider's perspective; D) it compares findings from a number of different cases
\\ \textit{Note:} The correct answer emphasizes that ethnographic research focuses on in-depth, immersive observation and engagement with participants, which allows researchers to gather nuanced insights and personal narratives that reflect the complexities of social life from the perspective of those being studied.
\item \textbf{Answer:} B
\\ \textit{Options:} A) pollies; B) mollies; C) dollies; D) lollies
\\ \textit{Note:} The term "mollies" refers specifically to a subculture of homosexual men in seventeenth and eighteenth century London, characterized by their distinct social gatherings and practices that deviated from the dominant heteronormative culture of the time.
\item \textbf{Answer:} D
\\ \textit{Options:} A) the product they were making; B) the production process; C) each other and humanity in general; D) all of the above
\\ \textit{Note:} Marx argued that capitalism alienates workers from the products of their labor, the labor process itself, their fellow workers, and their own human potential, leading to a comprehensive disconnection from all aspects of their existence.
\item \textbf{Answer:} D
\\ \textit{Options:} A) including union representatives in processes of policy decision making; B) creating links between the state and corporate organizations; C) recruiting women into full time paid employment; D) including working class organizations in political bargaining and representation
\\ \textit{Note:} The correct answer is supported by the understanding that 'incorporation' in the labor movement referred to the integration of working-class organizations into formal political processes, enabling them to participate in negotiations and influence policies that affected their rights and conditions.
\end{enumerate}

\section*{True / False}
\begin{enumerate}[label=\arabic*.]
\setcounter{enumi}{5}
\item \textbf{Answer:} False
\\ \textit{Options:} A) True; B) False
\\ \textit{Note:} Weber argued that sociologists should remain value-free in their research and analysis, emphasizing that knowledge is not inherently valuable but shaped by the context and perspectives of the researchers.
\item \textbf{Answer:} True
\\ \textit{Options:} A) True; B) False
\\ \textit{Note:} The human relations approach posits that fostering positive interpersonal relationships, effective communication, and prioritizing employee well-being significantly enhance motivation and productivity, ultimately contributing to organizational success.
\item \textbf{Answer:} False
\\ \textit{Options:} A) True; B) False
\\ \textit{Note:} The term 'culture industry,' coined by Theodor Adorno and Max Horkheimer, refers to the commodification and standardization of cultural products by capitalist mass media, rather than the emergence of diverse subcultures and counter-cultures.
\item \textbf{Answer:} True
\\ \textit{Options:} A) True; B) False
\\ \textit{Note:} The rise in cohabitation among couples reflects a shift in societal attitudes towards marriage, leading many to choose living together without formalizing their relationship through marriage, which contributes to the perception of declining marriage rates.
\item \textbf{Answer:} True
\\ \textit{Options:} A) True; B) False
\\ \textit{Note:} The statement is true because recent sociological studies indicate that drug use in Britain has expanded beyond traditional youth subcultures to include a broader demographic, reflecting changing social norms and patterns of consumption across various age groups and social classes.
\end{enumerate}

\section*{Long Form Response}
\begin{enumerate}[label=\arabic*.]
\setcounter{enumi}{10}
\item \textbf{Answer:} Socialization is the lifelong process through which individuals internalize the norms, values, and behaviors of their society, playing a crucial role in shaping personal identity and facilitating integration into the social fabric. This process is significantly influenced by various agents of socialization, including family, peers, schools, and media, each contributing unique perspectives and expectations that mold an individual's self-concept and social roles. For instance, family provides the foundational beliefs and values that guide early behavior, while peer groups offer a space for identity exploration and social acceptance. Educational institutions further reinforce societal norms and introduce individuals to a broader social context, while media shapes perceptions and aspirations. Collectively, these agents create a complex interplay that not only defines individual identities but also ensures cohesion within society, highlighting the profound impact of socialization on both personal and communal levels.
\\ \textit{Note:} correct because it comprehensively describes socialization as a lifelong process that integrates individuals into society and shapes their identity, while also identifying and elaborating on the specific agents of socialization that contribute to this process.
\item \textbf{Answer:} Relative deprivation, as proposed by Townsend (1979), refers to the perception of being disadvantaged in comparison to others within a society, which can lead to feelings of discontent and social unrest. One indicator of this phenomenon is media consumption, where individuals with lower socio-economic status may buy fewer than twenty DVDs in the previous year, reflecting their limited financial resources and access to cultural products. This disparity in media consumption not only highlights the material inequalities present in society but also suggests a broader cultural exclusion, as individuals unable to afford such items may feel disconnected from mainstream narratives and entertainment. Thus, Townsend's concept of relative deprivation underscores the importance of examining how access to cultural goods can both reflect and perpetuate social inequalities.
\\ \textit{Note:} The answer correctly identifies relative deprivation as a perception of disadvantage compared to others, illustrating how disparities in media consumption among different socio-economic groups serve as a reflection of social inequality and limited access to cultural resources.
\end{enumerate}

\vfill
\noindent\textit{End of Answer Key}
\end{document}